\section*{Exercise 1.2.3a}
\begin{proof}
    We first show green function for the first quadratic. As we need to put 2 negative charges for each boundary, a further positive charge is needed in the third quadratic. Next we show write out it's forms and check each properties. For $y \in \partial U, x\in U$,

\begin{equation}
    G(x,y) = \Phi(y-x) -  \Phi(y-x_2) - \Phi(y-x_4) + \Phi(y-x_3)
\end{equation}
where $x$ be the source point in the first quadrant, and $x_2, x_3, x_4$ be the corresponding symmetries of $x$ into the second, third, and fourth quadrants, respectively. Check the property: the harmonic property exists as we do not need to worry about the origin. Next, for $\Phi^{x}(y) = \Phi(y-x)$, if $y$ is on the x-axis, $\Phi(y-x_2) = \Phi(y-x_3)$, $\Phi(y-x) = \Phi(y-x_4)$. And if $y$ is on the y-axis, one has $\Phi(y-x_4) = \Phi(y-x_3)$, $\Phi(y-x_2) = \Phi(y-x)$. Hence, this meets our needs.

For much more general cases such that $n>2$, one may get $2^n-1$ additional chargers. Define $x_{(i)} = (i_1 x_1,\cdots, i_nx_n)$ where $i_j \in \{\pm 1\}$. One can always find a symmetric element to overcome the boudary condition. Then, we are done.
\end{proof}

For further steps to figure out it is a solution, as the poisson kernel's form is still similar to the form in Thm 1.15 ( upper plane problem), most of the steps should be the same as with $\int_{\partial U} K(x,y) dS(y) = 1$, and we will proof this equation in the next exercise.

\section*{Exercise 1.2.3b}
\begin{proof}
    We verify thing step by step. First regularity comes from harmonic, as \begin{equation*}
        \Delta_x K(x,y) = \partial_{\nu_y} \Delta_x G(x,y)
    \end{equation*} The exchangeability comes form boundedness of this region and continuity.Then one may exchange integral and differential sign again then one may verify that first two property hold.

    Before we start to proof third one, we first verify $\int_{\partial U} K(x,y) dS(y) = 1$ in general cases as $U$ open and bounded. As $K(x,y)$ have nothing to do with $g$, one may set $u = 1$ for any $x \in U$ and $u|_{\partial U} = 1$, we immediately see
    \begin{equation*}
        u(x) = 1 = \int_{\partial U} K(x,y) \cdot 1 dS(y)
    \end{equation*}
    Next, let $x_0\in \partial U, x \in U$ close to $x_0$, and some $\epsilon>0$, consider 
    \begin{align*}
        |u(x)-g(x_0)| & \leq \int_{\partial U} K(x,y) |\left(g(y)-g(x_0)\right)|dS(y)\\
                    &\leq \epsilon \int_{\partial U \cap B(x_0,\delta)} K(x,y) dS(y) +  2||g||_\infty\int_{\partial U \backslash B(x_0,\delta)} K(x,y) dS(y)
    \end{align*}
    As left part smaller than $\epsilon$ by above equality, what left is only the RHS. Still follow the argument on the note, we may let $|x-x_0| < \delta/2$ ,then one would have $|y-x| \geq |y-x_0|/2$, so that 
    \begin{align*}
      \frac{1-|x|^2}{n\alpha(n)} \int_{\partial U \backslash B(x_0,\delta)} \frac{1}{|y-x|^n}  dS(y) \leq \frac{2^n(1-|x|^2)}{n\alpha(n)} n \alpha(n) \frac{1}{\delta^n}
    \end{align*}
    Choosing some proper $\delta>0$, we are done here.
\end{proof}


\section*{Optional Question}
The question is similar with the ideas w.r.t. complex analysis. And one is not able to use cauchy type inequality.

Consider $\{r<|x|<R\}$, by Gauss-Green, one have for arbitrary $r<R$,
\begin{equation*}
    \int_{r<|x|<R} \Delta u dx = 0 = \int_{ \partial B(0,R)} \partial_{\nu} u dS -\int_{ \partial B(0,r)} \partial_{\nu} u dS
\end{equation*}
Hence, $F(R) = \int_{ \partial B(0,R)} \partial_{\nu} u dS$ is a constant. Define $A(r)$ as surface average,
$$ A(r) = \frac{1}{r^{n-1}n\alpha(n)} \int_{ \partial B(0,r)} u dS_r$$
$$A'(r) = \frac{1}{r^{n-1}n\alpha(n)} F(r) $$

So if $F \neq 0$, then one have $A$ behave like $\log r$ or $r^{2-n}$, which all not satisfy boundedness around origin. As $u\in C(\mathbb {R}^n)$, one have $\lim A(r) = u(0)$ as $r \to 0$. As $A$ is constant, one have Mean Value Property at 0, then we have function harmonic at 0, hence we are done by Liouville Theorem.


